\slajdid{04-s010}
\begin{frame}[label=04-s010]
\gusttekst
\setlength{\abovedisplayskip}{5pt}\setlength{\belowdisplayskip}{5pt}
Uočimo razlomljenu vrednost
\[\begin{aligned}
RV_{10}\cdot r&=\left(c_{-1}\cdot r^{-1}+c_{-2}\cdot r^{-2}+\ldots+c_{-k+1}\cdot r^{-k+1}+c_{-k}\cdot r^{-k}\right)\cdot r\\
&=c_{-1}\cdot r^0+c_{-2}\cdot r^{-1}+\ldots+c_{-k+1}\cdot r^{-k+2}+c_{-k}\cdot r^{-k+1}
\end{aligned}\]
\[c_{-1}\geq0\qquad c_i\in\{0,1,2,\ldots r-2,r-1\}\]
\[RV_{10}\cdot r=c_{-1}+(RV_{10})_{-2}=\text{celobrojna vrednost }c_{-1}+(RV_{10})_{-2}\]
Izdvojili smo cifru $c_{-1}$ i nastavljamo dalje
\[(RV_{10})_{-2}\cdot r=\text{celobrojna vrednost }c_{-2}+(RV_{10})_{-3}\]
Izdvojili smo cifru $c_{-2}$ i nastavljamo dalje \ldots

Zaustavljamo se kada bude zadovoljen uslov
\[(RV_{10})_{-k-1}=0\Rightarrow c_{-k-1}=0\]
\begin{center}Stalno bi dobijali prateće nule\end{center}
\end{frame}
